
\documentclass[10pt,a4paper,final, onecolumn]{article}
\usepackage[utf8]{inputenc}
\usepackage{amsmath}
\usepackage{amsfonts}
\usepackage{amssymb}
\usepackage{graphicx}
\usepackage{multicol}
\usepackage{quoting} 
\usepackage[LGR,T1]{fontenc}
\newcommand{\textgreek}[1]{\begingroup\fontencoding{LGR}\selectfont#1\endgroup}
\usepackage[a4paper,margin=0.75in]{geometry}
\usepackage{makecell}
\usepackage{tabularx}
\usepackage{hyperref}
\usepackage{dblfnote}


\renewcommand{\thefootnote}{\Alph{footnote}}
\renewcommand{\textsuperscript}[1]{\ensuremath{^{\textbf{\scriptsize #1}}}}



\begin{document}

The best-informed man is not necessarily the wisest. Indeed there is a danger that precisely in the multiplicity of his knowledge he will lose sight of what is essential. But on the other hand, knowledge of an apparently trivial detail quite often makes it possible to see into the depth of things. And so the wise man will seek to acquire the best possible knowledge about events, but always without becoming dependent upon this knowledge. To recognize the significant in the factual is wisdom. - Dietrich Bonhoeffer

\section*{Eccl 9:11-10:4 (NET)}


\subsection*{Wisdom Cannot Protect against Seemingly Chance Events}


    \textsuperscript{11} Again, I observed this on the earth\footnote{\textbf{tn} Heb “under the sun.”}: \\
    the race is not always\footnote{\textbf{tn} The term ``always” does not appear in the Hebrew text, but is supplied in the translation (five times in this verse) for clarity.} won by the swiftest, \\
    the battle is not always won by the strongest;\\\ 
    prosperity\footnote{\textbf{tn} Heb ``bread.”} does not always belong to those who are the wisest, \\
    wealth does not always belong to those who are the most discerning, \\
    nor does success\footnote{\textbf{tn} Heb ``favor.”} always come to those with the most knowledge— \\
    for time and chance may overcome them all. \\
    \textsuperscript{12} Surely, no one knows his appointed time!\footnote{\textbf{tn} Heb “time.” BDB 773 suggests that here (’et, ``time”) refers to an ``uncertain time.” On the other hand, HALOT 901 nuances it as ``destined time,” that is, ``no one knows his destined time [i.e., hour of destiny].”} \\
    Like fish that are caught in a deadly net, and like birds that are caught in a snare— \\
    just like them, all people are ensnared at an unfortunate time that falls upon them suddenly. \\


\subsection*{Most People Are Not Receptive to Wise Counsel}

    \textsuperscript{13} This is what I also observed about wisdom on earth, \\
    and it is a great burden to me: \\
    \textsuperscript{14} There was once a small city with a few men in it, \\
    and a mighty king attacked it, besieging it and building strong siege works against it. \\
    \textsuperscript{15} However, a poor but wise man lived in the city, \\
    and he could have delivered the city by his wisdom, \\
    but no one listened\footnote{\textbf{tn} Heb ``remembered.”} to that poor man. \\
    \textsuperscript{16} So I concluded that wisdom is better than might, \\
    but a poor man’s wisdom is despised; no one ever listens to his advice. \\


\subsection*{Wisdom versus Fools, Sin, and Folly}

    \textsuperscript{17} The words of the wise are heard in quiet, \\
    more than the shouting of a ruler is heard among fools. \\
    \textsuperscript{18} Wisdom is better than weapons of war, \\
    but one sinner can destroy much that is good. \\
    \textsuperscript{10:1} One dead fly makes the perfumer’s ointment give off a rancid stench,\footnote{\textbf{tn} The verb (ba’ash) means ``to cause to stink; to turn rancid; to emit a stinking odor” (e.g., Exod 16:24; Ps 38:6; Eccl 10:1)} \\
    so a little folly can outweigh much wisdom. \\


\subsection*{Wisdom Can Be Nullified By the Caprice of Rulers}

    \textsuperscript{2} A wise person’s good sense protects him,\footnote{\textbf{tn} Heb ``a wise man’s heart is at his right hand.” In ancient warfare, the shield of the warrior on one’s right-hand side protected one’s right hand. (e.g., Pss 16:8; 110:5; 121:5).} \\
    but a fool’s lack of sense leaves him vulnerable.\footnote{\textbf{tn} Heb ``and the heart of a fool is at his left hand.} \\
    \textsuperscript{3} Even when a fool walks along the road he lacks sense, \\
    and shows everyone what a fool he is. \\
    \textsuperscript{4} If the anger of the ruler flares up against you, do not resign from your position, \\
    for a calm response can undo great offenses. \\






%\end{paracol}
















\end{document}
